%%% FIELD def-ft-construction
For \(0<\lambda\leq\lambda_0\) and an i.i.d. observed-law sample
\(O_1,\ldots,O_N\), define
\[
\widehat\tau_{\lambda,N}=\frac1N\sum_{i=1}^N\psi_\lambda(O_i)
\]
using the public functions and the \(n\)-indexed certificates
\((r_{0n},r_{1n},r_{en})\).  Thus \(N\) is only the number of observations
averaged.  Set
\[
v_\kappa(\lambda)=
\begin{cases}
\lambda^{-(1-\kappa)/2},&0<\kappa<1,\\
\sqrt{1+\log(\lambda_0/\lambda)},&\kappa=1,\\
1,&\kappa>1,
\end{cases}
\qquad
\overline c_\kappa=\frac{c_+}{1-2^{-\kappa}},
\]
and
\[
H_\kappa=
\begin{cases}
\lambda_0^{-1}+\overline c_\kappa/(1-\kappa),&0<\kappa<1,\\
\lambda_0^{-1}+\overline c_1,&\kappa=1,\\
\lambda_0^{-1}+\overline c_\kappa\lambda_0^{\kappa-1}/(\kappa-1),&\kappa>1.
\end{cases}
\]
Define the finite-threshold envelope
\[
\begin{split}
U_N(\lambda)={}&
\frac{v_\kappa(\lambda)}{\sqrt N}
+\frac{r_{1n}}{\sqrt{N\lambda}}
+\frac{r_{en}}{\sqrt N\,\lambda^{3/2}}
+r_{1n}\lambda^{\kappa/2}
+\frac{r_{1n}r_{en}}{\lambda}
+r_{0n}r_{en}
\end{split}
\]
and the explicit score constant
\[
\begin{split}
C_{\mathrm{FT}}={}&1+\sqrt M+\sqrt{\sigma_1^2H_\kappa}
+\sqrt{\frac{\sigma_0^2}{\eta_+}}
+\frac{\bar r}{\sqrt{\eta_+}}
+\frac{2\bar r\sqrt{\sigma_0^2}}{\eta_+^{3/2}}\\
&+\sqrt2+\sqrt{\sigma_1^2}+\sqrt{\overline c_\kappa}
+\frac2{\eta_+}.
\end{split}
\]
For later oracle calibration, put \(v_0=v_\kappa(\lambda_0)\) and define
\[
C_{\mathrm{oe}}
=v_0+\bar r\lambda_0^{-1/2}+\bar r\lambda_0^{-3/2}
+2+\lambda_0^{-1}\bar r^{2/(\kappa+2)},
\]
\[
C_{\mathrm g}
=v_0+\lambda_0^{-1}+2\bar r\lambda_0^{-3/2}
+\lambda_0^{-(1-\kappa)/2}+1,
\]
\[
C_{\mathrm{opt}}=
\begin{cases}
4+C_{\mathrm{oe}}+C_{\mathrm g},&0<\kappa\leq1,\\
4+C_{\mathrm{oe}},&\kappa>1,
\end{cases}
\qquad
C_\star=C_{\mathrm{FT}}C_{\mathrm{opt}}.
\]

%%% FIELD lem-ft-statement
For every admissible radius triple
\((r_{0n},r_{1n},r_{en})\in[0,\bar r]^3\), let
\(\widetilde{\mathcal M}_n\) be the class obtained from
\texttt{def:public-weak-overlap-class} by deleting exactly
\texttt{ass:lower-shell-richness} and
\texttt{ass:upper-shell-richness}.  For every
\((P,m_0,m_1,p)\in\widetilde{\mathcal M}_n\), every integer \(N\geq1\),
every i.i.d. observed-law sample \(O_1,\ldots,O_N\), and every
\(0<\lambda\leq\lambda_0\), the block estimator and envelope from
\texttt{def:finite-threshold-envelope} satisfy the exact finite-threshold
inequality
\[
E_P\left|\widehat\tau_{\lambda,N}-\tau(P)\right|
\leq C_{\mathrm{FT}}U_N(\lambda).
\]
The public functions and all three approximation certificates in this display
remain indexed by \(n\); no membership assertion about any
sample-size-\(N\) model class is made.  Neither richness condition is used.

%%% FIELD lem-ft-proof
Fix a member of \(\widetilde{\mathcal M}_n\), an integer \(N\geq1\), and
\(0<\lambda\leq\lambda_0\).  All expectations below are under its observed
marginal.  First derive the two shell bounds required by the score calculation.
Write \(F(t)=P\{e(X)\leq t\}\).  By \texttt{ass:positivity}, the dyadic
shells \((2^{-j-1}t,2^{-j}t]\), \(j\geq0\), partition
\(\{0<e(X)\leq t\}\) up to a null set.  Therefore
\texttt{ass:uniform-shell} gives, for \(0<t\leq\lambda_0\),
\[
F(t)
\leq c_+t^\kappa\sum_{j=0}^\infty2^{-j\kappa}
=\overline c_\kappa t^\kappa.                                      \tag{1}
\]
Layer cake applied to the bounded nonnegative variable
\((e(X)\vee\lambda)^{-1}\), followed by the change of variables
\(s=t^{-1}\), gives
\[
E_P[(e(X)\vee\lambda)^{-1}]
=1+\int_\lambda^1\frac{P\{e(X)<t\}}{t^2}\,dt.                     \tag{2}
\]
On \([\lambda,\lambda_0]\) use (1), and on
\([\lambda_0,1]\) use the probability bound by one.  Evaluating the
resulting elementary integral separately for \(0<\kappa<1\), \(\kappa=1\),
and \(\kappa>1\) yields, with the constants defined in
\texttt{def:finite-threshold-envelope},
\[
E_P[(e(X)\vee\lambda)^{-1}]
\leq H_\kappa v_\kappa(\lambda)^2.                                  \tag{3}
\]
For example, when \(0<\kappa<1\), the right side before the final comparison
is at most
\(\lambda_0^{-1}+\overline c_\kappa\lambda^{\kappa-1}/(1-\kappa)\),
and \(v_\kappa(\lambda)^2=\lambda^{\kappa-1}\geq1\).  When
\(\kappa=1\), it is at most
\(\lambda_0^{-1}+\overline c_1\log(\lambda_0/\lambda)\).  When
\(\kappa>1\), it is at most
\(\lambda_0^{-1}+\overline c_\kappa\lambda_0^{\kappa-1}/(\kappa-1)\).
This proves (3) without either richness assumption.

Put
\[
\Delta=g_1-g_0,\qquad h_a=m_a-g_a,
\qquad \varepsilon_a=Y-g_a(X)\quad\hbox{on }\{D=a\},
\qquad e_\lambda=e\vee\lambda.
\]
The definitions of \(g_a\) give
\(E_P(\varepsilon_a\mid D=a,X)=0\).  Since
\texttt{ass:right-overlap} places \(e(X)\) in
\([0,1-\eta_+]\), projection onto \([0,1-\eta_+/2]\) and the map
\(z\mapsto z\vee\lambda\) are contractions at \(e(X)\).  Hence
\[
|p^\circ-e|\leq|p-e|,
\qquad |q_\lambda-e_\lambda|\leq|p-e|,
\qquad 1-p^\circ\geq\eta_+/2.                                      \tag{4}
\]
All denominators below are positive by \(\lambda>0\),
\texttt{ass:positivity}, and \texttt{ass:right-overlap}.

Substitution in \texttt{def:clipped-score}, followed by adding and subtracting
the true clipped treated denominator and the true control denominator, gives
the pointwise identity
\[
\psi_\lambda-\tau(P)=A+B_1+B_0+C_1+C_0+R_1+R_0,                    \tag{5}
\]
where
\[
\begin{aligned}
A&=\Delta-\tau(P)+\frac{D\varepsilon_1}{e_\lambda}
-\frac{(1-D)\varepsilon_0}{1-e},\\
B_1&=h_1\left(1-\frac D{e_\lambda}\right),
&B_0&=-h_0\left(1-\frac{1-D}{1-e}\right),\\
C_1&=D\varepsilon_1\left(\frac1{q_\lambda}-\frac1{e_\lambda}\right),
&C_0&=-(1-D)\varepsilon_0
\left(\frac1{1-p^\circ}-\frac1{1-e}\right),\\
R_1&=Dh_1\left(\frac1{e_\lambda}-\frac1{q_\lambda}\right),
&R_0&=(1-D)h_0
\left(\frac1{1-p^\circ}-\frac1{1-e}\right).
\end{aligned}                                                       \tag{6}
\]
Conditional mean zero of the residuals shows
\(E_PA=E_PB_0=E_PC_1=E_PC_0=0\).  The only clipping mean is
\[
E_PB_1
=E_P\!\left[h_1(X)\left\{1-\frac{e(X)}{e_\lambda(X)}\right\}\right].
\]
By Cauchy--Schwarz, \texttt{ass:public-g1-radius}, and (1),
\[
|E_PB_1|
\leq r_{1n}F(\lambda)^{1/2}
\leq\sqrt{\overline c_\kappa}\,r_{1n}\lambda^{\kappa/2}.          \tag{7}
\]

We next give every second-moment calculation used for the centered empirical
terms.  By Minkowski, \texttt{ass:bounded-treatment-effect-signal}, the two
conditional residual-variance assumptions, \texttt{ass:right-overlap}, and
(3),
\[
\begin{split}
\|A\|_{P,2}
&\leq\sqrt M
+\sqrt{\sigma_1^2E_P[e_\lambda^{-1}]}
+\sqrt{\sigma_0^2/\eta_+}\\
&\leq
\left(\sqrt M+\sqrt{\sigma_1^2H_\kappa}
+\sqrt{\sigma_0^2/\eta_+}\right)v_\kappa(\lambda),
\end{split} \tag{8}
\]
where \(v_\kappa(\lambda)\geq1\).  Conditional on \(X\),
\[
E_P\!\left[\left.h_1^2\left(1-\frac D{e_\lambda}\right)^2
\right|X\right]
\leq h_1^2(1+\lambda^{-1}),
\]
and
\[
E_P\!\left[\left.h_0^2\left(1-\frac{1-D}{1-e}\right)^2
\right|X\right]
=h_0^2\frac e{1-e}\leq h_0^2\eta_+^{-1}.
\]
The public regression certificates and \(\lambda\leq\lambda_0<1\) thus imply
\[
\|B_1\|_{P,2}\leq\sqrt2\,r_{1n}\lambda^{-1/2},
\qquad
\|B_0\|_{P,2}\leq r_{0n}\eta_+^{-1/2}.                            \tag{9}
\]
Using (4), \(e\leq e_\lambda\), the public propensity certificate, and the
conditional residual-variance bounds gives
\[
\begin{split}
E_PC_1^2
&\leq\sigma_1^2E_P\!\left[
\frac{e(p-e)^2}{q_\lambda^2e_\lambda^2}\right]
\leq\sigma_1^2r_{en}^2\lambda^{-3},\\
E_PC_0^2
&\leq\sigma_0^2E_P\!\left[
\frac{(p^\circ-e)^2}{(1-p^\circ)^2(1-e)}\right]
\leq4\sigma_0^2r_{en}^2\eta_+^{-3}.
\end{split} \tag{10}
\]
Finally, conditioning on \(X\), using (4), and applying Cauchy--Schwarz to the
two certified \(L^2(P_X)\) errors yields the first-moment remainder bounds
\[
E_P|R_1|\leq\frac{r_{1n}r_{en}}\lambda,
\qquad
E_P|R_0|\leq\frac2{\eta_+}r_{0n}r_{en}.                             \tag{11}
\]

For every mean-zero square-integrable random variable \(Z\), independence of
the \(N\) observations and Cauchy--Schwarz give
\[
E_P\left|\frac1N\sum_{i=1}^NZ_i\right|
\leq\left\{E_P\left[\left(\frac1N\sum_{i=1}^NZ_i\right)^2\right]\right\}^{1/2}
=\frac{\|Z\|_{P,2}}{\sqrt N}.                                     \tag{12}
\]
Apply (12) to \(A,B_0,C_1,C_0\) and to \(B_1-E_PB_1\), use
\(\|B_1-E_PB_1\|_{P,2}\leq\|B_1\|_{P,2}\), retain (7), and use
\(E_P|N^{-1}\sum_iR_{a,i}|\leq E_P|R_a|\).  Since
\(r_{0n},r_{en}\leq\bar r\) and \(v_\kappa(\lambda)\geq1\),
(8)--(12) give
\[
\begin{split}
E_P|\widehat\tau_{\lambda,N}-\tau(P)|
\leq{}&K_v\frac{v_\kappa(\lambda)}{\sqrt N}
+\sqrt2\frac{r_{1n}}{\sqrt{N\lambda}}
+\sqrt{\sigma_1^2}\frac{r_{en}}{\sqrt N\lambda^{3/2}}\\
&+\sqrt{\overline c_\kappa}r_{1n}\lambda^{\kappa/2}
+\frac{r_{1n}r_{en}}\lambda
+\frac2{\eta_+}r_{0n}r_{en},
\end{split} \tag{13}
\]
where
\[
K_v=\sqrt M+\sqrt{\sigma_1^2H_\kappa}
+\sqrt{\sigma_0^2/\eta_+}
+\bar r\eta_+^{-1/2}
+2\bar r\sqrt{\sigma_0^2}\eta_+^{-3/2}.
\]
By its displayed algebraic definition, \(C_{\mathrm{FT}}\) is at least each
of \(K_v,\sqrt2,\sqrt{\sigma_1^2},
\sqrt{\overline c_\kappa},1,2/\eta_+\).  Comparing (13) term by term with
\(U_N(\lambda)\) proves
\[
E_P|\widehat\tau_{\lambda,N}-\tau(P)|
\leq C_{\mathrm{FT}}U_N(\lambda),
\]
as claimed.  Every condition used above is retained in
\(\widetilde{\mathcal M}_n\), and neither richness inequality entered any
step.

%%% FIELD proof-upper
Let \(\widetilde{\mathcal M}_n\) be the broader class in the statement.
The causal target is the frozen identified functional
\(\tau(P)=E_P[g_1(X)-g_0(X)]\); the statistic below is a sample map and is not
a definition of that target.  Its ATE interpretation follows from
\texttt{ass:consistency} and \texttt{ass:ignorability}.

Apply \texttt{lem:finite-threshold-clipped-bound} with \(N=n\).  For every
\(0<\lambda\leq\lambda_0\),
\[
\sup_{(P,m_0,m_1,p)\in\widetilde{\mathcal M}_n}
E_P|\widehat\tau_{\lambda,n}-\tau(P)|
\leq C_{\mathrm{FT}}U_n(\lambda).                                  \tag{14}
\]
It remains only to evaluate the explicit envelope at the frozen oracle
threshold.  Put
\[
b_o=n^{-1/(\kappa+1)},\qquad
b_e=r_{en}^{2/(\kappa+2)},\qquad
b_g=(nr_{1n}^2)^{-1},
\]
with \(b_g=+\infty\) when \(r_{1n}=0\), and let \(L\) be the active maximum
in \texttt{def:oracle-threshold}.  Also abbreviate
\[
S_o=r_{1n}n^{-\kappa/[2(\kappa+1)]},\qquad
S_e=r_{1n}r_{en}^{\kappa/(\kappa+2)},\qquad
S_0=r_{0n}r_{en}.
\]
Then \(\rho_n=Q_\kappa(n,r_{1n})+S_o+S_e+S_0\).

Suppose first that \(L<\lambda_0\), so
\(\lambda_n^\star=L\).  Since \(L\geq b_o\),
\[
\frac{r_{1n}}{\sqrt{nL}}
\leq r_{1n}L^{\kappa/2}.                                           \tag{15}
\]
Since \(L\geq b_e\), with the assertion trivial when \(r_{en}=0\),
\[
\frac{r_{en}}{\sqrt nL^{3/2}}
\leq\frac{L^{(\kappa-1)/2}}{\sqrt n}
\leq Q_\kappa(n,r_{1n}),
\qquad
\frac{r_{1n}r_{en}}L\leq S_e.                                     \tag{16}
\]
For the second inequality in the first part of (16), if \(0<\kappa<1\),
use \(L\geq b_g\) to obtain
\(n^{-1/2}L^{-(1-\kappa)/2}
\leq n^{-\kappa/2}r_{1n}^{1-\kappa}\); if \(\kappa\geq1\), use
\(L<\lambda_0<1\).  The same comparison with \(b_g\) gives
\[
\frac{v_\kappa(L)}{\sqrt n}\leq Q_\kappa(n,r_{1n})                 \tag{17}
\]
for \(0<\kappa\leq1\).  At \(\kappa=1\), indeed,
\(\log(\lambda_0/L)\leq\log_+(nr_{1n}^2)\).  For \(\kappa>1\),
(17) is equality because both sides equal \(n^{-1/2}\).

Taking the positive \(\kappa/2\)-power of the active maximum and bounding a
maximum by a sum yields
\[
r_{1n}L^{\kappa/2}
\leq S_o+S_e+Q_\kappa(n,r_{1n}).                                    \tag{18}
\]
When \(0<\kappa<1\), the contribution of \(b_g\) in (18) is exactly
\(n^{-\kappa/2}r_{1n}^{1-\kappa}\); when \(\kappa=1\), it is
\(n^{-1/2}\); and when \(\kappa>1\), \(b_g\) is not active.  Combining
(15)--(18) term by term in \(U_n(L)\) gives the explicit non-cap bound
\[
U_n(L)\leq4Q_\kappa(n,r_{1n})+2S_o+3S_e+S_0
\leq4\rho_n.                                                       \tag{19}
\]

Now suppose the cap binds, so \(L\geq\lambda_0\) and
\(\lambda_n^\star=\lambda_0\).  At least one active component is at least
\(\lambda_0\).  If \(b_o\geq\lambda_0\), then the fourth term of
\(U_n(\lambda_0)\) is at most \(S_o\).  If
\(b_e\geq\lambda_0\), that term is at most \(S_e\).  In either case,
using \(Q_\kappa(n,r_{1n})\geq n^{-1/2}\), both radius caps, and
\[
\frac{r_{1n}r_{en}}{\lambda_0}
\leq\lambda_0^{-1}\bar r^{2/(\kappa+2)}S_e,
\]
we obtain directly from the six terms of \(U_n\)
\[
U_n(\lambda_0)\leq C_{\mathrm{oe}}\rho_n.                           \tag{20}
\]
If \(0<\kappa\leq1\) and \(b_g\geq\lambda_0\), then
\(r_{1n}\leq(n\lambda_0)^{-1/2}\), including the case \(r_{1n}=0\).
Consequently
\[
\begin{gathered}
\frac{r_{1n}}{\sqrt{n\lambda_0}}
\leq\lambda_0^{-1}n^{-1/2},\qquad
r_{1n}\lambda_0^{\kappa/2}
\leq\lambda_0^{-(1-\kappa)/2}n^{-1/2},\\
\frac{r_{1n}r_{en}}{\lambda_0}
\leq\bar r\lambda_0^{-3/2}n^{-1/2}.
\end{gathered}
\]
Together with the first and third stochastic terms and \(S_0\), these
inequalities give
\[
U_n(\lambda_0)\leq C_{\mathrm g}\rho_n.                            \tag{21}
\]
The cases (20)--(21) exhaust the active maximum.  Equations (19)--(21) and the
definition of \(C_{\mathrm{opt}}\) prove the fully explicit optimization
inequality
\[
U_n(\lambda_n^\star)\leq C_{\mathrm{opt}}\rho_n.                    \tag{22}
\]

Taking \(\lambda=\lambda_n^\star\) in (14), using (22), and recalling
\(C_\star=C_{\mathrm{FT}}C_{\mathrm{opt}}\), gives
\[
\sup_{(P,m_0,m_1,p)\in\widetilde{\mathcal M}_n}
E_P|\widehat\tau_{\lambda_n^\star}-\tau(P)|
\leq C_\star\rho_n.                                               \tag{23}
\]
Thus the unnamed constant \(C\) in the theorem statement may be taken to be
the computable algebraic constant \(C_\star\) from
\texttt{def:finite-threshold-envelope}.  The statistic is measurable in the
observed sample, public functions, and certified radii, so
\texttt{def:estimator-class} gives
\(\widehat\tau_{\lambda_n^\star}\in\mathcal T_n\).

For \(W=|\widehat\tau_{\lambda_n^\star}-\tau(P)|\), Markov's inequality and
(23) yield, for every fixed \(\delta\in(0,1/4)\),
\[
P\{W>C_\star\rho_n/\delta\}
\leq\frac{\delta E_PW}{C_\star\rho_n}\leq\delta.                   \tag{24}
\]
Finally, \(\mathcal M_n\subseteq\widetilde{\mathcal M}_n\) because the former
adds exactly the two richness conditions and changes no retained member
condition.  Restricting the suprema in (23)--(24) proves both conclusions on
\(\mathcal M_n\).  Neither richness condition was used for the upper bound.

%%% FIELD proof-sharp
All upper and lower bounds use the frozen law--function class
\(\mathcal M_n\) and decision class \(\mathcal T_n\).  By
\texttt{thm:clipped-absolute-upper}, with its explicit witness
\(C_\star\),
\[
\mathfrak R_n^{(1)}
\leq\sup_{(P,m_0,m_1,p)\in\mathcal M_n}
E_P|\widehat\tau_{\lambda_n^\star}-\tau(P)|
\leq C_\star\rho_n.                                               \tag{25}
\]

For the root-MSE upper bound, suppose first that \(n\geq5\), and let \(N_j\)
be the size of block \(j\) in \texttt{def:median-estimator}.  Since the five
block sizes differ by at most one,
\[
N_j\geq\lfloor n/5\rfloor\geq n/10.                               \tag{26}
\]
The last inequality follows directly for \(5\leq n\leq9\), while for
\(n\geq10\) it follows from
\(\lfloor n/5\rfloor\geq n/5-1\geq n/10\).  In the explicit envelope
\(U_N(\lambda)\), exactly the first three terms depend on \(N\), and each is
proportional to \(N^{-1/2}\); the last three are deterministic biases.
Therefore (26), \texttt{lem:finite-threshold-clipped-bound}, and \texttt{lem:oracle-envelope-optimization} give, for every block,
\[
\begin{split}
E_P|\widehat\tau_{\lambda_n^\star,N_j}-\tau(P)|
&\leq C_{\mathrm{FT}}U_{N_j}(\lambda_n^\star)\\
&\leq\sqrt{10}\,C_{\mathrm{FT}}U_n(\lambda_n^\star)\\
&\leq\sqrt{10}\,C_\star\rho_n.
\end{split} \tag{27}
\]
Here only the averaging size changes to \(N_j\); the public functions,
certificates, threshold, and rate remain indexed by the original \(n\).
The blocks are independent by \texttt{ass:iid}.  Apply the already established
\texttt{lem:median-risk-upgrade} with
\(K=\sqrt{10}C_\star\) to obtain
\[
\sup_{(P,m_0,m_1,p)\in\mathcal M_n}
\left\{E_P[(\widehat\tau_{\mathrm{med}}-\tau(P))^2]\right\}^{1/2}
\leq\sqrt{210}\,C_\star\rho_n.                                   \tag{28}
\]

For \(1\leq n<5\), the median rule is the full-sample clipped rule.  Put
\[
\lambda_{\min}=\min\{\lambda_0,4^{-1/(\kappa+1)}\}>0.
\]
Since \(b_o\geq4^{-1/(\kappa+1)}\), the oracle threshold satisfies
\(\lambda_n^\star\geq\lambda_{\min}\), while
\(1-p^\circ\geq\eta_+/2\).  With
\(\Delta=g_1-g_0\), \(h_a=m_a-g_a\), and
\(\varepsilon_a=Y-g_a(X)\) on \(\{D=a\}\), direct substitution gives
\[
\psi_{\lambda_n^\star}
=\Delta+h_1-h_0
+\frac{D(\varepsilon_1-h_1)}{q_{\lambda_n^\star}}
-\frac{(1-D)(\varepsilon_0-h_0)}{1-p^\circ}.                       \tag{29}
\]
The signal moment, the two public regression certificates, the two conditional
residual-variance bounds, and Minkowski's inequality imply
\[
\|\psi_{\lambda_n^\star}-\tau(P)\|_{P,2}\leq C_{\mathrm{sm}},      \tag{30}
\]
where the following displayed finite constant suffices:
\[
C_{\mathrm{sm}}
=2\sqrt M+2\bar r
+\lambda_{\min}^{-1}\sqrt{2(\sigma_1^2+\bar r^2)}
+\frac2{\eta_+}\sqrt{2(\sigma_0^2+\bar r^2)}.
\]
Indeed, \(|\tau(P)|\leq\sqrt M\), and the two weighted residual-regression
terms in (29) have \(L^2(P)\) norms bounded by the last two summands before
centering.  Pointwise convexity of the square now gives
\[
E_P\!\left[\left(\frac1n\sum_{i=1}^n
\{\psi_{\lambda_n^\star}(O_i)-\tau(P)\}\right)^2\right]
\leq\frac1n\sum_{i=1}^n
E_P[(\psi_{\lambda_n^\star}(O_i)-\tau(P))^2]
\leq C_{\mathrm{sm}}^2.                                           \tag{31}
\]
Because \(\rho_n\geq n^{-1/2}\geq1/2\) for \(n<5\), (31) is at most
\((2C_{\mathrm{sm}}\rho_n)^2\).  Combining (28)--(31), membership of the
median rule from \texttt{lem:median-risk-upgrade}, and the square-inside-the-
expectation definition \texttt{def:root-mse-risk}, proves
\[
\mathfrak R_n^{(2)}\leq C_2\rho_n                                  \tag{32}
\]
for all \(n\geq1\), with \(C_2\) depending only on the fixed class constants.

For fixed \(\delta\in(0,1/4)\), (24) makes
\(C_\star\rho_n/\delta\) feasible in
\texttt{def:probability-risk}; hence
\[
\mathfrak R_{n,\delta}^{(\mathrm{prob})}
\leq(C_\star/\delta)\rho_n.                                       \tag{33}
\]

For the converse, let \(c>0\) and \(\delta_*>0\) be the uniform constants in
\texttt{thm:four-channel-converse}, and set
\(\delta_0=\min\{\delta_*/2,1/8\}\).  That theorem directly gives
\[
\mathfrak R_n^{(1)}\geq c_1\rho_n,
\qquad
\mathfrak R_n^{(2)}\geq c_2\rho_n                                 \tag{34}
\]
for positive fixed constants.  Moreover, for every \(T\in\mathcal T_n\),
\[
\sup_{(P,m_0,m_1,p)\in\mathcal M_n}
P\{|T-\tau(P)|\geq c\rho_n\}
\geq\delta_*\geq2\delta_0.                                       \tag{35}
\]
If \(0<\delta\leq\delta_0\) and \(u<c\rho_n\), the event in (35) is
contained in \(\{|T-\tau(P)|>u\}\), whose supremum probability is therefore
strictly larger than \(\delta\).  Such a radius is infeasible in
\texttt{def:probability-risk}; taking the two infima gives
\[
\mathfrak R_{n,\delta}^{(\mathrm{prob})}\geq c\rho_n.               \tag{36}
\]
Equations (25), (32)--(34), and (36) prove all three equivalences.  The upper
construction, including the blockwise use of (27), remains valid on the
broader uniform one-sided-shell class.  The matching lower bounds use the
reservoir-rich class and both richness assumptions through
\texttt{thm:four-channel-converse}.

%%% FIELD proof-honest
Let \(\widetilde{\mathcal M}_n\) be the broader class in the statement.  We
first prove coverage there using the explicit interval calibration in the
corrected \texttt{def:bias-aware-interval}.

Suppose \(n\geq5\), and let \(N_j\) be a block size.  Because the five block sizes differ by at most one, \(N_j\geq\lfloor n/5\rfloor\geq n/10\): the second inequality is immediate for \(5\leq n\leq9\), and for \(n\geq10\) it follows from \(\lfloor n/5\rfloor\geq n/5-1\geq n/10\).  Only the first three terms of \(U_N(\lambda)\) scale with \(N^{-1/2}\), while the last three are unchanged biases.  Hence the declared lemmas \texttt{lem:finite-threshold-clipped-bound} and \texttt{lem:oracle-envelope-optimization} give \[\begin{split}
E_P|\widehat\tau_{\lambda_n^\star,N_j}-\tau(P)|
&\leq C_{\mathrm{FT}}U_{N_j}(\lambda_n^\star)\\
&\leq\sqrt{10}\,C_{\mathrm{FT}}U_n(\lambda_n^\star)\\
&\leq\sqrt{10}\,C_\star\rho_n.
\end{split} \tag{37}
\]
The five block estimators are independent by \texttt{ass:iid}.  Applying
\texttt{lem:median-risk-upgrade} with
\(K=\sqrt{10}C_\star\) yields, for every \(u>0\),
\[
\sup_{(P,m_0,m_1,p)\in\widetilde{\mathcal M}_n}
P\{|\widehat\tau_{\mathrm{med}}-\tau(P)|>u\}
\leq\min\!\left\{1,10\left(\frac{\sqrt{10}C_\star\rho_n}{u}\right)^3\right\}.
                                                                        \tag{38}
\]
For \(1\leq n<5\), the median rule is the full-sample clipped rule, and the explicit mean bound in \texttt{thm:clipped-absolute-upper} plus Markov gives
\[
\sup_{(P,m_0,m_1,p)\in\widetilde{\mathcal M}_n}
P\{|\widehat\tau_{\mathrm{med}}-\tau(P)|>u\}
\leq C_\star\rho_n/u.                                             \tag{39}
\]

The corrected definition sets
\[
K_\alpha=C_\star\max\!\left\{
\sqrt{10}(10/\alpha)^{1/3},\alpha^{-1}\right\}.
\]
Taking \(u=K_\alpha\rho_n\) in (38) makes its second argument at most
\(\alpha\), and taking the same \(u\) in (39) also gives a bound by
\(\alpha\).  Therefore, for every \(n\geq1\),
\[
\inf_{(P,m_0,m_1,p)\in\widetilde{\mathcal M}_n}
P\{\tau(P)\in\mathcal C_{n,\alpha}\}\geq1-\alpha.                 \tag{40}
\]
The center is measurable in the sample, and the radius is an explicit function
of the supplied radii and fixed class constants.  Thus
\(\mathcal C_{n,\alpha}\in\mathcal I_n\), and its deterministic length is
\[
\operatorname{length}(\mathcal C_{n,\alpha})
=2K_\alpha\rho_n.                                                  \tag{41}
\]
Equations (40)--(41) prove the upper coverage and expected-length conclusions
on the broader class without either richness condition.

We now restrict to the nonempty reservoir-rich class \(\mathcal M_n\).  The
four component experiments certified in the proof of
\texttt{thm:four-channel-converse} are two-point or fuzzy two-prior
experiments generated by \(\mathfrak H_{\mathrm{lb}}\), and every fixed-sign
law belongs to \(\mathcal M_n\).  For a two-point component, let its targets
be separated by \(2d\) and let the total variation distance between its two
sample laws be at most \(\Gamma\).  Honesty at both laws and total-variation
transfer imply, under the first law,
\[
P_0\{\tau(P_0)\in C_n\}\geq1-\alpha,
\qquad
P_0\{\tau(P_1)\in C_n\}\geq1-\alpha-\Gamma.
\]
An interval containing both targets has length at least \(2d\), so the union
bound gives
\[
P_0\{\operatorname{length}(C_n)\geq2d\}
\geq1-2\alpha-\Gamma.                                              \tag{42}
\]

For a fuzzy component, let its two sample mixtures have total variation at
most \(\Gamma\), and let their prior target centers be separated by \(2d\).
The finite cell count in \(\mathfrak H_{\mathrm{lb}}\) can be chosen so that,
under either prior, the fixed-sign target is within \(d/4\) of its center
except on a set of prior probability at most \(\varepsilon\).  Averaging
fixed-law honesty under each prior, transferring one event between mixtures,
and intersecting the two events show that under the first mixture the interval
meets both \(d/4\)-neighborhoods with probability at least
\(1-2\alpha-\Gamma-2\varepsilon\).  Their distance is at least \(3d/2\);
hence
\[
\overline P_-
\{\operatorname{length}(C_n)\geq3d/2\}
\geq1-2\alpha-\Gamma-2\varepsilon.                                 \tag{43}
\]
The established converse construction chooses the two-point amplitudes with
\(\Gamma\leq1/4\), and in each fuzzy experiment enlarges its finite cell count
until \(\Gamma+2\varepsilon\leq1/4\), without changing target separation or
fixed-sign membership.  Since \(\alpha<1/4\), the right sides of (42)--(43)
are bounded below by a positive constant depending only on \(\alpha\).  A
mixture expectation is an average of fixed-law expectations, so (43) implies
that at least one fixed-sign member has expected interval length bounded below
by a positive constant times \(d\).

The four target separations in \texttt{thm:four-channel-converse} dominate,
up to fixed positive constants, the four summands of \(\rho_n\).  If their
scales are \(d_1,\ldots,d_4\), then
\[
\max_{1\leq j\leq4}d_j
\geq\frac14\sum_{j=1}^4d_j
\geq c\rho_n.                                                       \tag{44}
\]
Selecting the largest component in (44) and applying (42) or (43) proves that
every interval rule honest over \(\mathcal M_n\) satisfies
\[
\sup_{(P,m_0,m_1,p)\in\mathcal M_n}
E_P[\operatorname{length}(C_n)]\geq c_\alpha\rho_n.                 \tag{45}
\]
Taking the infimum over honest rules and combining (45) with (41), restricted
through \(\mathcal M_n\subseteq\widetilde{\mathcal M}_n\), gives
\(\mathfrak L_{n,\alpha}\asymp\rho_n\).  Only the lower bound and equivalence
use the two richness assumptions.

%%% FIELD def-ci-change
For \((r_{0n},r_{1n},r_{en})\in[0,\bar r]^3\), use the explicit constant
\(C_\star\) from \texttt{def:finite-threshold-envelope}, set
\[
K_\alpha=C_\star\max\left\{
\sqrt{10}(10/\alpha)^{1/3},\alpha^{-1}\right\},
\qquad
\mathcal C_{n,\alpha}
=\left[\widehat\tau_{\mathrm{med}}-K_\alpha\rho_n,
\widehat\tau_{\mathrm{med}}+K_\alpha\rho_n\right],
\]
and define
\[
\mathfrak L_{n,\alpha}
=\inf_{C_n\in\mathcal I_n:\,
\inf_{(P,m_0,m_1,p)\in\mathcal M_n}P\{\tau(P)\in C_n\}\geq1-\alpha}
\sup_{(P,m_0,m_1,p)\in\mathcal M_n}
E_P[\operatorname{length}(C_n)].
\]

%%% FIELD def-ci-replacement
For \((r_{0n},r_{1n},r_{en})\in[0,\bar r]^3\), use the explicit constant
\(C_\star\) from \texttt{def:finite-threshold-envelope}, set
\[
K_\alpha=C_\star\max\left\{
\sqrt{10}(10/\alpha)^{1/3},\alpha^{-1}\right\},
\qquad
\mathcal C_{n,\alpha}
=\left[\widehat\tau_{\mathrm{med}}-K_\alpha\rho_n,
\widehat\tau_{\mathrm{med}}+K_\alpha\rho_n\right],
\]
and define
\[
\mathfrak L_{n,\alpha}
=\inf_{C_n\in\mathcal I_n:\,
\inf_{(P,m_0,m_1,p)\in\mathcal M_n}P\{\tau(P)\in C_n\}\geq1-\alpha}
\sup_{(P,m_0,m_1,p)\in\mathcal M_n}
E_P[\operatorname{length}(C_n)].
\]

%%% FIELD stmt-upper
For every admissible radius triple \((r_{0n},r_{1n},r_{en})\in[0,\bar r]^3\), let \(\widetilde{\mathcal M}_n\) denote the broader class obtained from \texttt{def:public-weak-overlap-class} by retaining all of its member conditions and deleting exactly \texttt{ass:lower-shell-richness} and \texttt{ass:upper-shell-richness}.  The rule \(\widehat\tau_{\lambda_n^{\star}}\) from \texttt{def:clipped-score} and \texttt{def:oracle-threshold} belongs to \(\mathcal T_n\).  There is a constant \(C<\infty\), depending only on the retained fixed class constants, including \(M\) and \(\bar r\), such that
\[
\sup_{(P,m_0,m_1,p)\in\widetilde{\mathcal M}_n}
E_P\!\left|\widehat\tau_{\lambda_n^{\star}}-\tau(P)\right|
\leq C\rho_n.
\]
For every fixed \(\delta\in(0,1/4)\), the same estimator satisfies
\[
\sup_{(P,m_0,m_1,p)\in\widetilde{\mathcal M}_n}
P\!\left\{\left|\widehat\tau_{\lambda_n^{\star}}-\tau(P)\right|
>C\rho_n/\delta\right\}\leq\delta.
\]
Since \(\mathcal M_n\subseteq\widetilde{\mathcal M}_n\), both risk conclusions remain valid with \(\widetilde{\mathcal M}_n\) replaced by the reservoir-rich class \(\mathcal M_n\).

%%% FIELD stmt-sharp
Uniformly over the bounded radius regime \((r_{0n},r_{1n},r_{en})\in[0,\bar r]^3\) for which \(\mathcal M_n\) is nonempty and over the fixed reservoir-rich shell constants, with every minimax infimum taken over \(T\in\mathcal T_n\) measurable in the observed sample and allowed to use \(m_0\), \(m_1\), \(p\), and the certified radii, \(\mathfrak R_n^{(1)}\asymp\rho_n\), \(\mathfrak R_n^{(2)}\asymp\rho_n\), and \(\mathfrak R_{n,\delta}^{(\mathrm{prob})}\asymp\rho_n\) for every fixed \(\delta\in(0,\delta_0]\). Thus the same \(\rho_n\) is the sharp frontier for absolute risk, root mean-squared risk, and fixed-probability error.

%%% FIELD stmt-honest
For every admissible radius triple \((r_{0n},r_{1n},r_{en})\in[0,\bar r]^3\), let \(\widetilde{\mathcal M}_n\) denote the broader class obtained from \texttt{def:public-weak-overlap-class} by retaining all of its member conditions and deleting exactly \texttt{ass:lower-shell-richness} and \texttt{ass:upper-shell-richness}.  For every fixed \(\alpha\in(0,1/4)\), a choice of \(K_\alpha\) depending only on \(\alpha\), \(\bar r\), and the retained fixed class constants makes \(\mathcal C_{n,\alpha}\in\mathcal I_n\) uniformly honest on this broader class:
\[
\inf_{(P,m_0,m_1,p)\in\widetilde{\mathcal M}_n}
P\{\tau(P)\in\mathcal C_{n,\alpha}\}\geq1-\alpha,
\qquad
\sup_{(P,m_0,m_1,p)\in\widetilde{\mathcal M}_n}
E_P[\operatorname{length}(\mathcal C_{n,\alpha})]
\leq C_\alpha\rho_n.
\]
Conversely, when the reservoir-rich class \(\mathcal M_n\) is nonempty, every \(C_n\in\mathcal I_n\) that is uniformly honest over \(\mathcal M_n\), including every rule using the public functions and certified radii, satisfies
\[
\sup_{(P,m_0,m_1,p)\in\mathcal M_n}
E_P[\operatorname{length}(C_n)]\geq c_\alpha\rho_n.
\]
Consequently \(\mathfrak L_{n,\alpha}\asymp\rho_n\) on the reservoir-rich class.  No matching lower bound or length equivalence is asserted here for \(\widetilde{\mathcal M}_n\setminus\mathcal M_n\).
